\begin{abstract}

Polynomial systems are a fundamental language for describing algebraic and geometric problems across mathematics and its applications.
Beyond finding individual solutions, one often seeks global properties of the solution space,
such as its dimension and irreducible components,
or structural properties and decompositions of the underlying polynomial ideal.
% [solving large sets of polynomial equations is important for problems in X Y Z]
% Solving a large set of polynomial equations is a fundamental problem in algebraic geometry and a pratical problem in many applications.
% Computing a Gr\"obner basis of a polynomial ideal is a fundamental problem in algebraic geometry.
% [many of those problems are still inaccessible \ldots]
For large polynomial systems with many variables and high degrees,
determining these global properties can become intractable.
% [leading approaches include X Y Z]
% [benchmark relative to other methods]
Leading approaches include variants of Buchberger's algorithm with syzygy, decomposition techniques, resultants, and XXX,
as well as numerical methods such as homotopy continuation.
% [however, these approaches are limited by \ldots]
However, even the most modern approaches still often fail on practical problems and on slight perturbations of standard benchmarks.
% [broadest strongest claim] and [we demonstrate it on X Y]
Here we introduce a framework that enables exact algebraic analysis of polynomial systems beyond the reach of existing methods.
%Here we introduce a framework for Gr\"obner-basis computation and use it
We use this framework to compute bases for two systems arising from research in quantum field theory
for which, no other standard method has succeeded.
% [more specifics about the method]
The framework divides the computation into smaller subproblems that existing methods can solve.
It also draws on conflict-driven clause learning (CDCL) to reduce the number of subproblems.
% [outlook sentence on what would this contribute to in the figure / extensions for future work]
At a later stage, we construct a decision tree guided by domain-informed hypotheses to efficiently enumerate candidate solutions.
%Several heuristics in the framework remain to be optimized, and further ideas from CDCL may also apply.
Together, these ideas offer a path toward the exact analysis of previously inaccessible polynomial systems.
\end{abstract}

% update to include mentions of other important achievements including proofs if nontrivial
